% ==============================================================================
% SECAO 1 — INTRODUCAO
% ==============================================================================
\section{INTRODUÇÃO}\label{sec:intro}

\subsection{Contextualização e Motivação}\label{sec:intro_contexto}
A expansão acelerada de aplicações conectadas no ecossistema da Internet das Coisas (IoT), abrangendo redes de sensores industriais, telemetria veicular, monitoramento energético e dispositivos urbanos inteligentes, consolidou as plataformas de processamento contínuo de dados em \textit{streaming} como infraestrutura crítica da computação moderna. Diferentemente dos paradigmas tradicionais em lote, esses sistemas recebem fluxos ininterruptos de eventos que precisam ser ingeridos, armazenados temporariamente em memória e processados em tempo quase real pelos operadores computacionais, frequentemente sob acordos de nível de serviço (\textit{Service Level Agreements}, SLAs).

Na operação prática dessas arquiteturas, a taxa de chegada de eventos raramente permanece constante ao longo do tempo, estando sujeita a oscilações estocásticas, surtos de tráfego e degraus persistentes de demanda. Quando a capacidade de processamento alocada é inferior à taxa de entrada, os eventos excedentes acumulam-se nas filas temporárias (\textit{buffers}) de ingestão. Esse acúmulo prolongado eleva de forma direta a latência fim a fim, degrada a experiência das aplicações consumidoras, compromete o cumprimento dos SLAs e, em cenários extremos, conduz à exaustão de memória e ao descarte de pacotes. Em contrapartida, a adoção de superprovisionamento estático, que mantém continuamente recursos ociosos para absorver picos eventuais, resulta em desperdício de energia elétrica e de custos financeiros de nuvem.

Nesse cenário, a Teoria de Controle Realimentado oferece um arcabouço matemático consolidado para viabilizar o escalonamento dinâmico automático (\textit{autoscaling}) de recursos computacionais (DORF; BISHOP, 2018; FRANKLIN; POWELL; EMAMI-NAEINI, 2013). Ao monitorar continuamente o comprimento da fila e realimentar essa informação para ajustar a capacidade solicitada à infraestrutura, torna-se possível estabilizar o nível de congestionamento em patamares seguros mesmo na presença de perturbações permanentes de carga. Contudo, a síntese desse controlador exige considerar as características físicas do ambiente computacional, em especial a natureza integradora do acúmulo em fila, a inércia temporal associada ao provisionamento de novos operadores e os limites rígidos de saturação impostos pela capacidade instalada.

\subsection{Problema Analisado, Objetivos e Organização do Trabalho}\label{sec:intro_objetivos}
Este relatório investiga a modelagem dinâmica, a análise de estabilidade e o projeto de controladores clássicos em malha fechada para uma plataforma de processamento de eventos IoT em \textit{streaming}. O sistema opera em torno de um ponto nominal caracterizado por taxa de chegada e capacidade de processamento equilibradas em $\lambda_0 = \mu_0 = 1000\text{ eventos/s}$, com fila inicial acumulada de $q(0) = 600\text{ eventos}$. O mecanismo de ajuste de recursos responde aos comandos de alocação com dinâmica de primeira ordem e constante de tempo $\tau = 2\text{ s}$, estando fisicamente limitado ao intervalo operacional $\mu(t) \in [600,\,1400]\text{ eventos/s}$. Deseja-se conduzir e manter a fila no patamar de referência $q_r = 300\text{ eventos}$ (correspondente a um incremento de redução $x_r = 300\text{ eventos}$), garantindo desempenho transitório rápido e bem amortecido mesmo após a incidência de uma perturbação permanente de carga $\Delta\lambda = +100\text{ eventos/s}$ em $t = 20\text{ s}$.

Como objetivos específicos, o estudo propõe-se a: (i) deduzir o modelo físico-matemático incremental da planta e sua representação em diagrama de blocos; (ii) caracterizar a resposta natural em malha aberta e demonstrar a inviabilidade de regulação sem realimentação; (iii) avaliar o desempenho do controlador puramente proporcional (P) mediante o critério de Routh-Hurwitz e o Lugar Geométrico das Raízes (LGR), detalhando seus compromissos estruturais; (iv) projetar analiticamente e comparar no plano complexo e no domínio da frequência (Bode) compensadores P, PI, PD e PID; (v) investigar computacionalmente os fenômenos não lineares impostos pela saturação de atuadores, \textit{integral windup} e \textit{derivative kick}; e (vi) identificar a arquitetura de controle com melhor compromisso global de engenharia frente às metas operacionais.

O restante do trabalho está organizado da seguinte forma. A Seção~\ref{sec:modelagem_fisica} apresenta a arquitetura física da plataforma e deduz as equações de conservação de fluxo e de alocação de capacidade. A Seção~\ref{sec:malha_aberta} unifica o modelo em uma equação diferencial de segunda ordem e analisa o comportamento da planta em malha aberta. Na Seção~\ref{sec:controle_p}, investiga-se o controle proporcional sob as perspectivas de Routh-Hurwitz, resposta modal e LGR. A Seção~\ref{sec:projeto_comp} formaliza as metas quantitativas de desempenho, detalha a sintonia analítica dos compensadores P, PI, PD e PID, examina suas margens frequenciais e demonstra suas limitações estruturais e não lineares. A Seção~\ref{sec:simulacao} apresenta as simulações temporais completas ao longo de $60\text{ s}$ e consolida a tabela comparativa de métricas. Por fim, a Seção~\ref{sec:conclusoes} sintetiza os compromissos de engenharia envolvidos e apresenta as considerações finais do estudo.
